\section{Related Work}
\label{sec:related}

\subsection{Linear representations, sparse features, and their skeptics}

A growing body of evidence supports the linear-representation view that high-level concepts are encoded as directions in the residual stream $x \in \mathbb{R}^d$, recoverable by linear probes and difference-of-means contrasts \citep{marks2023geometry}. Sparse autoencoders (SAEs) extend this picture by decomposing activations into an overcomplete dictionary of approximately monosemantic features \citep{bricken2023monosemanticity,cunningham2023sparse,templeton2024scaling}, and open suites now make such dictionaries available across layers and architectures \citep{lieberum2024gemmascope,gao2024scaling,rajamanoharan2024jumprelu}. Within this line, several latents have been linked to epistemic and knowledge-awareness behavior relevant to hallucination, for instance entity-recognition features that fire on known versus unknown entities \citep{ferrando2025entity}. Our hallucination/epistemic subspace $H_t$, contrasting answerable real-entity against surface-matched unanswerable fictional-entity questions, targets the same phenomenon, though we recover it by contrast at each dialogue turn rather than reading off a fixed dictionary. A parallel skeptical literature cautions that the interpretability value of these features is easy to overstate: SAE-derived probes can underperform simple baselines on downstream tasks \citep{kantamneni2025useful}, and common SAE proxy metrics do not reliably track the downstream quantities they are meant to predict \citep{karvonen2025saebench}. This tension, between features that look meaningful in isolation and features that earn their keep under controlled evaluation, motivates the measurement framing we adopt: a geometric signal that appears present must be shown to survive the right null before it is treated as real.

\subsection{Steering and concept erasure}

A complementary line treats these directions as handles for intervention rather than objects of study. Representation engineering reads and writes concept directions to control model behavior \citep{zou2023repe}, with task-specific instances including inference-time intervention for truthfulness \citep{li2023iti}, contrastive activation addition \citep{panickssery2024caa}, and the identification of a single refusal direction whose ablation via weight orthogonalization removes safety behavior \citep{arditi2024refusal}. The dual operation, removing a concept, is formalized by closed-form concept erasure \citep{belrose2023leace} and nullspace projection \citep{ravfogel2020inlp}. These methods presuppose that the relevant concept occupies a stable, low-dimensional subspace that an intervention can target. Our safety subspace $S_t$ (harmful versus surface-matched harmless requests) and the static disentangling projector we stress-test are built directly on this presupposition: the projector is fit once and asked to hold across turns. Whether such a static target is adequate, or whether it must be re-estimated as dialogue content changes, is precisely the question a conditional treatment of these subspaces raises.

\subsection{Conditional steering and conditional disentanglement}

The premise we test belongs to a recent and still thin line that makes intervention \emph{conditional} on dialogue content: the claim that the safety and hallucination geometry shifts with what has been said, so that a context-dependent (rather than static) subspace is needed to steer or to disentangle reliably. The closest method is \citet{mahmoud2025unintended}, which combines SAEs with subspace orthogonalization to disentangle hallucination from safety, with the goal of intervening on one without unintended effects on the other. This is the method whose premise our study evaluates: if the relevant subspaces drift with content, a static disentangler should leak as context accumulates, motivating a conditional one. We emphasize that this is a single recent preprint, so the conditional-disentanglement hypothesis is best regarded as an open conjecture rather than an established regularity, and our contribution is to ask whether its motivating premise holds under confound control on the models we evaluate, rather than to propose a competing intervention.

\subsection{Brittleness of steering and the geometry of refusal}

Several results already complicate any simple static-versus-conditional dichotomy, from both directions. On one side, steering vectors are reported to be brittle and to generalize poorly out of distribution \citep{tan2024steering}, which could be read as evidence that a static direction is insufficient and that conditioning on context is required. On the other side, the refusal direction is not as singular as the one-direction picture suggests: refusal appears to be multi-directional \citep{pan2025hidden}, and concept representations form cones in which geometric orthogonality between two concepts does not imply that interventions on them are independent \citep{wollschlager2025geometry}. These findings matter for our measurement because they decouple two things that a naive reading conflates. Apparent failure of a static subspace, or apparent movement of one, need not be content-conditionality: it can instead reflect multi-directional structure, cone geometry, or the gauge freedom in choosing an orthonormal basis $U_t$. We therefore work with gauge-invariant quantities throughout, the chordal drift $d(U,V) = \sqrt{\operatorname{mean}_i \sin^2 \theta_i}$ and the cross-overlap $\Omega(S_t,H_t) = \lVert U_S^\top U_H \rVert_F^2 / k$ defined via principal angles \citep{edelman1998geometry}, so that a reported drift cannot be an artifact of how the bases were oriented. The lesson we draw from this literature is that brittleness and apparent drift are necessary but not sufficient evidence for content-conditionality, and that distinguishing a genuine content effect from a confound requires an explicit null.

\subsection{Measurement, confound control, and our contribution}

Where prior work largely assumes a representational structure and then intervenes on it, we treat the structure itself as the object of measurement and ask whether a specific structural claim, content-conditionality of $H_t$ and $S_t$, survives controls. The central methodological point is the choice of null. The naive null, two bootstrap resamples of the same probe set at a fixed context, carries zero between-condition length and topic variance, so on a between-context statistic the comparison ``spread exceeds floor'' is near-tautological. We instead require that genuine content-conditionality exceed a scrambled-twin null that word-shuffles each user turn, preserving the exact token multiset (length, topic, and lexicon) while destroying coherent meaning, with dialogue-level bootstrap confidence intervals so the verdict does not rest on a single dialogue. This protocol isolates three confounds, the no-context probe-in-context onset, token length and final-token position, and single-dialogue idiosyncrasy, and we show that each, left uncontrolled, manufactures a positive result that the next control overturns. Validated against analytic ground truth (a recovered $15^\circ$/turn injected rotation, and a static subspace correctly returning drift at the bootstrap floor) and powered through a reported noise floor, the protocol returns \texttt{NOT\_CONTENT\_CONDITIONAL} on every setting we evaluate: Qwen2.5-0.5B-Instruct and Qwen2.5-1.5B-Instruct \citep{qwen25} in raw activations, and GPT-2 small in the SAE feature space of residual SAEs, the space these methods actually use. A static disentangling projector that appears to fail across turns leaks no more than its length and topic scrambled twin, so its apparent failure is a confound rather than evidence that a content-conditional projector is needed. Our scope is deliberately bounded to these instruction-tuned and GPT-2 models and to rank-1 difference-of-means subspaces; against benchmarks such as TruthfulQA \citep{lin2022truthfulqa} and over-refusal probes like XSTest \citep{rottger2024xstest} that drive much of the steering literature, the contribution here is not a new direction or a new intervention but a reusable, pre-registered protocol and a set of controls that let future work tell a real content effect apart from a length and topic artifact before building conditional machinery on top of it.